\par\Needspace{6\baselineskip}\medskip
\begin{center}\hypertarget{native-196-article146}{}\pdfbookmark[2]{146.}{article-146}146.\end{center}
For the calculation of the elements in the first combination we use $r',r'',2f$ and the corrected time interval between the second and third observations. We denote the product of this interval with the quantity $k$ (Art. 1) by $\theta$. In the second combination we use $r,r',2f''$ and the interval between the first and second observations, setting its product with $k$ $=\theta''$. In each case the method set out in Arts.$^{\prime}$ 88—105 need be followed only as far as the quantity denoted there by $y$. We denote its value in the former combination by $\eta$ and in the latter by $\eta''$. Then set
\[
\frac{\theta''\eta}{\theta\eta''}=P',\qquad
\frac{r'r'\theta\theta''}{rr''\eta\eta''\cos f\cos f'\cos f''}=Q',
\]
\SourceLeaf{197}It is clear that if the values of the quantities $P,Q$ on which the entire calculation thus far was based were themselves correct, we must have \mbox{$P'=P,Q'=Q$}. Conversely, it is easy to see that if we obtain $P'=P,Q'=Q$, the two calculations of the elements, when both are carried through to the end, will yield exactly equal numbers. These will therefore represent all three observations exactly and fully satisfy the problem. Whenever $P'=P,Q'=Q$ does not hold, take $P'-P,Q'-Q$ for $X$ and $Y$ if $P$ and $Q$ were taken for $x$ and $y$. It will be more convenient still to set $\log P=x$, $\log Q=y$, $\log P'-\log P=X$, $\log Q'-\log Q=Y$.
The calculation must then be repeated with other values of $x,y$.
